\chapter{GIT e GitHub}
\label{chapter:git-e-github}
In questo capitolo vengono spiegati due strumenti quali \tb{Git} e \tb{GitHub}, molto utili per la sincronizzazione e condivisione online di file, pur non essendo strettamente legati al tema Linux sono comunque utili da conoscere in quanto molto supportati da questo sistema, anche dal punto di vista filosofico, semplici, efficaci e versatili.

\section{GIT}
\label{sec:git}
% breve intro sulla storia e cosa fa in breve
Il controllo di versione è una procedura che consiste nel tracciare il modo con cui i dati in un file vengono modificati nel tempo, GIT è molto forte in questo, con il controllo attivo è possibile ripristinare in qualsiasi momento lo stato del file in corrispondenza della data specificata. Il metodo con cui GIT tiene traccia delle versioni dei file è innovativo e risolve un problema molto importante ovvero quello di dover avere file di backup 
, talvolta se il file deve poter essere modificato da più utenti, tutti i backup devono trovarsi in un server e questo rende il tutto ancora più complicato da gestire; GIT si dice essere un sistema distribuito ovvero il processo di \ti{file versioning} è organizzato in modo tale che ogni utente riesca a modificare in locale il file, registrare le modifiche attraverso GIT e poi, quando è possibile connettersi ad un server, più nello specifico GitHub, il file modificato e le versioni vengono aggiornate lasciando ad altri utenti il file aggiornato, anche dal punto di vista del \ti{versioning}

 ma non solo, infatti l'innovativo modo con cui GIT effettua il \ti{versioning} permette di avere completo controllo anche su grandi progetti dove molte persone collaborano e quindi eventi come modifiche simultanee ai file possono essere comuni e non banali da gestire (nel contesto di \ti{file versioning}), per questo GIT si dice essere distribuito, in pratica ogni utente che possiede il progetto possiede anche il versionamento completo di questo . 


Allo stesso tempo GIT è pensato anche per permettere la condivisione dei file con versione in modo che questi possano essere modificati  


 è stato sviluppato avendo sempre la possibilità di ritornare, in qualsiasi momento, ad un punto 
\begin{itemize}
	\item \tb{Controllo di versione}:
	\item \tb{Distribuito}:
	\item \tb{Controllo di versione}:
	\item \tb{Controllo di versione}:
	\item \tb{Controllo di versione}: 
\end{itemize}
